\subsection{Модуль lab5}

\subsubsection{EulerianCycle: проверка и нахождение эйлерова цикла}

\textbf{Вход (конструктор):} \texttt{Graph\&}~--- неориентированный граф.

\textbf{Выход (\texttt{findCycle}):} \texttt{std::optional<vector<int>>}~---
последовательность вершин эйлерова цикла или пути;
\texttt{std::nullopt} если граф не является ни эйлеровым, ни полуэйлеровым.

\textbf{Описание:} Метод \texttt{isEulerian} проверяет отсутствие
вершин нечётной степени и связность графа; \texttt{isSemiEulerian}~---
ровно две вершины нечётной степени при том же условии связности.
Связность определяется BFS-обходом по ненулевым вершинам через
\texttt{hasSingleNonZeroComponent}.

Метод \texttt{makeEulerian} выполняет достройку: сначала соединяет
несвязные компоненты цепочкой рёбер, затем итеративно устраняет
вершины нечётной степени попарным добавлением рёбер. При выборе пар
предпочтение отдаётся вершинам, между которыми ещё нет ребра. Все
добавленные рёбра фиксируются в \texttt{m\_added\_edges\_}.

Метод \texttt{hierholzer} реализует алгоритм Иерхольцера итеративно
на стеке с вектором \texttt{usedEdge} и указателями \texttt{ptr[u]}
для пропуска использованных рёбер без повторного перебора.
Сложность: $O(V + E)$, пространственная: $O(V + E)$.

\begin{lstlisting}[style=cppstyle, caption={Достройка и алгоритм Иерхольцера (EulerianCycle.cc)}]
bool EulerianCycle::isEulerian() const;
bool EulerianCycle::isSemiEulerian() const;
void EulerianCycle::makeEulerian();
std::optional<std::vector<int>> EulerianCycle::hierholzer(
    int start, int totalEdges) const;

stk.push(id2idx[start]);
while (!stk.empty()) {
    int u = stk.top();
    bool found = false;
    while (ptr[u] < static_cast<int>(g[u].size())) {
        auto [v, eid] = g[u][ptr[u]++];
        if (usedEdge[eid]) continue;
        usedEdge[eid] = true;
        stk.push(v);
        found = true;
        break;
    }
    if (!found) {
        circuit.push_back(vertices[u]);
        stk.pop();
    }
}
std::reverse(circuit.begin(), circuit.end());
\end{lstlisting}

\subsubsection{CutSystem: фундаментальная система разрезов}

\textbf{Вход (конструктор):} \texttt{Graph const\& originalGraph},
\texttt{Graph const\& mst}~--- исходный граф и его минимальный остов.

\textbf{Выход (\texttt{buildFundamentalCuts}):}
\texttt{vector<FundamentalCut>}~--- вектор из $|V|-1$ структур с
полями \texttt{treeEdge}, \texttt{leftComponent},
\texttt{rightComponent} и \texttt{cutEdges}.

\textbf{Описание:} Для каждого ребра $e \in T$ метод
\texttt{componentWithoutTreeEdge} BFS-обходом по остову без ребра $e$
находит компоненту $A$; дополнение $B = V \setminus A$ строится
перебором всех вершин. Рёбра исходного графа, пересекающие границу
$A \mid B$, собираются в \texttt{cutEdges}, нормализуются и
дедуплицируются. Сложность: $O(|V| \cdot (|V| + |E|))$.

Метод \texttt{symmetricDifference} кодирует каждое ребро как
\texttt{uint64} и поддерживает множество \texttt{active} по принципу
XOR: при повторном вхождении код удаляется, при первом~--- добавляется.
Сложность: $O(k \cdot |\mathrm{Cut}|)$.

\begin{lstlisting}[style=cppstyle, caption={Построение ФСР и симметрическая разность (CutSystem.cc)}]
std::vector<CutSystem::FundamentalCut>
    CutSystem::buildFundamentalCuts() const;
std::vector<std::pair<int, int>> CutSystem::symmetricDifference(
    std::vector<FundamentalCut> const& cuts,
    std::vector<int> const& selectedIndices) const;

for (auto const& edge : mst_.edges()) {
    auto treeEdge = normalizeEdge(edge.from, edge.to);
    auto left = componentWithoutTreeEdge(treeEdge.first, treeEdge);
    for (auto const& e : originalGraph_.edges()) {
        auto ge = normalizeEdge(e.from, e.to);
        if (left.contains(ge.first) != left.contains(ge.second))
            cut.cutEdges.push_back(ge);
    }
}

for (int idx : selectedIndices)
    for (auto const& edge : cuts[idx].cutEdges) {
        auto code = encode(edge);
        if (active.contains(code)) active.erase(code);
        else                       active.insert(code);
    }
\end{lstlisting}
